% Generated draft: requires mathematical and source review.
\documentclass[11pt]{article}
\usepackage[margin=25mm]{geometry}
\usepackage{fontspec}
\setmainfont{Latin Modern Roman}
\usepackage{amsmath,amssymb,mathtools,unicode-math}
\setmathfont{Latin Modern Math}
\usepackage{enumitem,needspace,xurl,hyperref}
\setlength{\parindent}{0pt}
\setlength{\parskip}{5pt}
\setlength{\emergencystretch}{3em}
\newcommand{\N}{\mathbb{N}}
\newcommand{\Z}{\mathbb{Z}}
\newcommand{\Q}{\mathbb{Q}}
\newcommand{\R}{\mathbb{R}}
\newcommand{\C}{\mathbb{C}}
\newcommand{\F}{\mathbb{F}}
\newcommand{\A}{\mathbb{A}}
\newcommand{\PP}{\mathbb{P}}
\newcommand{\E}{\mathbb{E}}
\newcommand{\eps}{\varepsilon}
\newcommand{\dd}{\,\mathrm{d}}
\begin{document}
\section*{UnsolvedMath ID: 30005497}
\textbf{Source problem code:} catalogue entry
\subsection{Definitions and mathematical statement}
An indecomposable principally polarized abelian variety is a principally polarized abelian variety \((A,\Theta)\) that is not isomorphic, as a polarized variety, to a product of two positive-dimensional principally polarized abelian varieties. Here \(\Theta\) is an effective divisor representing the principal polarization, and write \(g=\dim A\).

The supplied record does not specify the ground field or define \(\operatorname{Sing}_m(\Theta)\). For a complex formulation, the expected multiplicity-locus interpretation is
\[
\operatorname{Sing}_m(\Theta)=\{x\in\Theta:\operatorname{mult}_x(\Theta)\geq m\},
\]
where \(\operatorname{mult}_x(\Theta)\) is the order of vanishing of a local defining equation of \(\Theta\) at \(x\). These interpretations require confirmation against the report.

With these qualifications, the supplied question is: for every indecomposable principally polarized abelian variety \((A,\Theta)\), does
\[
\exists\,m\in\Z,\quad m\geq2,\qquad
\dim\operatorname{Sing}_m(\Theta)=g-2m+1
\]
imply that \(A\) is the Jacobian of a hyperelliptic curve or the intermediate Jacobian of a cubic threefold?

In the usual complex interpretation, the proposed models are \(J(C)=\operatorname{Pic}^0(C)\) for a smooth connected projective hyperelliptic curve \(C\), and
\[
J(X)=\frac{H^3(X,\C)}{F^2H^3(X,\C)+\operatorname{im}(H^3(X,\Z)\to H^3(X,\C))}
\]
for a smooth cubic threefold \(X\subset\PP^4_{\C}\), with their canonical principal polarizations. The supplied wording does not explicitly specify smoothness or whether the conclusion requires an isomorphism preserving the principal polarization.

There is also an unresolved boundary qualification: the record gives no convention for \(\dim\varnothing\), no nonemptiness hypothesis on \(\operatorname{Sing}_m(\Theta)\), and no upper bound on \(m\). Since \(g-2m+1\) may be negative, the displayed premise requires clarification in the empty-locus case. No additional restriction on \(m\) is imposed here.

The original excerpt refers to equality in an omitted inequality numbered \((1)\). It also states that an affirmative answer would imply a conjecture of Debarre, proposed by Grushevsky in \([8,\text{ Conjecture }5.5]\), and a conjecture of Grushevsky in \([8,\text{ Conjecture }5.12]\). Neither the omitted inequality nor the bibliographic identity of reference \([8]\) is supplied.
\subsection{Short English statement}
The question asks whether equality \(\dim\operatorname{Sing}_m(\Theta)=\dim A-2m+1\) for some integer \(m\geq2\) forces an indecomposable principally polarized abelian variety to be a hyperelliptic Jacobian or a cubic-threefold intermediate Jacobian. The ground field, multiplicity-locus notation, empty-locus boundary case, and precise meaning of the proposed identifications need confirmation.
\subsection{Sources}
\begin{itemize}
\item Supplied local UnsolvedMath registry record, permanent ID \(30005497\), title: Equality Cases for Singularities of Theta Divisors; source problem code \texttt{catalogue entry}; collection: Oberwolfach Reports Open Problems; proposed year: 2023. No separate proposer is recorded.
\item Supplied authoritative report citation: \emph{Komplexe Analysis -- Differential and Algebraic methods in Kähler spaces} (2023), Oberwolfach Reports, pp.\ 965--1030, DOI \texttt{10.4171/owr/2023/18}. \url{https://doi.org/10.4171/owr/2023/18}. The excerpt is attributed to Ruijie Yang, joint work with Christian Schnell, under the title \emph{The Riemann--Schottky problem via singularities of theta divisors}.
\item Supplied literature source labeled ``Schnell and Yang'': \url{https://arxiv.org/abs/2309.16763}.
\item Supplied literature source labeled ``OWR 2023 report'': \url{https://ems.press/content/serial-article-files/47013}.
\item Literature assessment dated \texttt{2026-08-22}: the supplied curation labels the problem partially solved and states that the dimension bound and equality problem are proved in Prym, low-dimensional, and isolated-singularity cases, while remaining open in general. This is supplied dated curation, not an independently verified literature assessment; the precise scope and hypotheses of those cases require review.
\item The registry labels the statement \texttt{corrected\_verified} and reports that it was checked against the official Oberwolfach report. This is a supplied verification claim, not independent verification performed for this draft. The original excerpt's reference \([8]\), including Conjectures 5.5 and 5.12, lacks a supplied bibliographic entry.
\end{itemize}
\end{document}
